\documentclass{report}
\usepackage{amsmath,amssymb}
\setlength{\parindent}{0mm}
\setlength{\parskip}{1em}
\begin{document}
\begin{center}
\rule{6in}{1pt} \
{\large Grey Ballard \\
{\bf Reducing Communication Costs for Sparse Matrix Multiplication within Algebraic Multigrid}}

P O Box 969 \\ Mail Stop 9159 \\ Livermore \\ CA 94550
\\
{\tt gmballa@sandia.gov}\\
Grey Ballard\\
Jonathan Hu\\
Christopher Siefert\end{center}

We consider the sequence of sparse matrix-matrix multiplications
performed during the setup phase of algebraic multigrid.
In particular, we show that the most commonly used parallel algorithm is
often not the most communication-efficient one for all of the matrix
multiplications involved.
By using an alternative algorithm, we show that the communication costs
are reduced (in theory and practice), and we demonstrate the performance
benefit for both model and real problems on large-scale
distributed-memory parallel systems.

We consider the smoothed aggregation multigrid method, and for this talk
we focus on a symmetric fine grid operator $A$.
After fine-level rows of $A$ are grouped into coarse-level aggregates,
there are three sparse matrix multiplication operations performed to
construct the coarse grid operator.
The tentative prolongation matrix $\widehat{P}$ represents aggregate
membership, and the final prolongation matrix $P$ is computed by applying
a step of damped Jacobi to $\widehat{P}$, which is structurally
equivalent to computing the product $A\cdot \widehat{P}$.
The coarse grid operator is given by the triple product $A_c=P^TAP$,
which is typically performed with two more sparse matrix multiplications:
$A\cdot P$ and $P^T\cdot (AP)$.

The standard approach for performing each of these parallel sparse matrix
multiplications is to use a row-wise algorithm: for general $C=A\cdot B$,
each processor owns a subset of the rows of $A$, a subset of the rows of
$B$, and computes a subset of rows of $C$ (which matches the distribution
of $A$).
The communication required is an exchange of rows of $B$ among processors.
We consider as an alternative the outer-product sparse matrix
multiplication algorithm, where each processor owns a subset of the rows
of $A$, the corresponding subset of columns of $B$, and the communication
consists of exchanging partially unreduced rows of $C$ among processors
(assuming the desired output distribution is row-wise).

Because of its low arithmetic intensity, the performance of parallel
algorithms for sparse matrix multiplication is typically bound by the
cost of interprocessor communication.
Based on communication cost analysis of model problems, we conclude that
the row-wise algorithm is the right choice for computing $A\cdot
\widehat{P}$ and $A\cdot P$ but the outer-product is the better choice
for computing $P^T\cdot (AP)$.
For a fine grid operator corresponding to a 3D 27-point stencil on a
regular grid, for example, the row-wise algorithm for $P^T(AP)$ requires
more than $5\times$ the communication of the outer-product algorithm.
Furthermore, using the outer-product algorithm for $P^T\cdot(AP)$ avoids
the explicit redistribution of $P^T$ that is typically required in order
to use the row-wise algorithm and involves as much interprocessor
communication as a sparse matrix multiplication.

We have implemented the outer-product algorithm within the Trilinos
software framework, and we compare communication costs and running times
of the two approaches for representative matrices.
To confirm the theoretical analysis, our experiments include regular-grid
stencil matrices for 2D and 3D problems.
We also perform tests on realistic problems, including a low Mach fluid
dynamics application problem, to demonstrate the benefits of the new
approach.


\end{document}
